\section{Three Doors, One Column}
\label{sec:ch12-three-doors-one-column}

\index{authorization!the guard on the root field|(}%
The speaker list is not public. Every speaker in this conference has an
address the program committee writes to, it is nobody's business but the
committee's, and the Speakers service is where it lives. So \code{Speaker}
gains a column and the service gains a way to tell one caller from another.

\index{authorization!two packages}%
Two packages, one for each half of that. \code{Speakers.csproj} again, with
what moved marked:

\begin{minted}[fontsize=\footnotesize,highlightlines={13,18}]{xml}
<Project Sdk="Microsoft.NET.Sdk.Web">

  <PropertyGroup>
    <TargetFramework>net10.0</TargetFramework>
    <Nullable>enable</Nullable>
    <ImplicitUsings>enable</ImplicitUsings>
    <GenerateDocumentationFile>true</GenerateDocumentationFile>
    <NoWarn>$(NoWarn);CS1591</NoWarn>
  </PropertyGroup>

  <ItemGroup>
    <PackageReference Include="HotChocolate.ApolloFederation" Version="16.6.1" />
    <PackageReference Include="HotChocolate.AspNetCore" Version="16.6.1" />
    <PackageReference Include="HotChocolate.AspNetCore.Authorization" Version="16.6.1" />
    <PackageReference Include="HotChocolate.AspNetCore.CommandLine" Version="16.6.1" />
    <PackageReference Include="HotChocolate.Data.EntityFramework" Version="16.6.1" />
    <PackageReference Include="HotChocolate.Types.Analyzers" Version="16.6.1" />
    <PackageReference Include="Microsoft.AspNetCore.Authentication.JwtBearer" Version="10.0.11" />
    <PackageReference Include="Microsoft.EntityFrameworkCore.Sqlite" Version="10.0.11" />
  </ItemGroup>

</Project>
\end{minted}

\index{authorization!the development signing key}%
The token is a JWT signed with a key both this service and, later, the router
have to hold. A deployment reads that key out of a secret store. This book
stands one up nowhere, so the key is a constant in a file of its own, which
is also the honest place for it: a secret with a name like this one is not
pretending to be a secret.

\begin{minted}{csharp}
namespace Speakers;

/// <summary>
/// The symmetric key this service validates tokens with. A real
/// deployment reads a key from configuration and validates issuer,
/// audience and lifetime; this one is fixed so that a reader can mint
/// a token and reproduce every request in the chapter.
/// </summary>
public static class SigningKey
{
    public const string Value =
        "splitting-the-graph-development-signing-key-32b";
}
\end{minted}

\index{Program.cs@\code{Program.cs}!wiring authentication in}%
Then \code{Program.cs}, complete, with the four blocks this chapter adds:

\begin{minted}[highlightlines={1,5,28-46,49,52-54,75-76}]{csharp}
using System.Text;
using ChilliCream.Nitro.App;
using HotChocolate.AspNetCore;
using Microsoft.EntityFrameworkCore;
using Microsoft.IdentityModel.Tokens;
using Speakers;

var builder = WebApplication.CreateBuilder(args);

// EF Core reports every command through the logging pipeline as well,
// at Information and with a duration attached. StatementLog below is
// this service's instrument, so silence the other one rather than have
// every statement printed twice.
builder.Logging.AddFilter(
    "Microsoft.EntityFrameworkCore.Database.Command",
    LogLevel.Warning);

builder.Services.AddDbContextFactory<SpeakerContext>(options =>
{
    options.UseSqlite("Data Source=speakers.db");

    if (builder.Environment.IsDevelopment())
    {
        options.AddInterceptors(new StatementLog());
    }
});

// A shared secret rather than a JWKS endpoint, because this book
// stands up no identity provider and a reader has to be able to mint
// a token and repeat every request in the chapter. The three
// Validate flags are off for the same reason and are the three a
// deployment turns on first. What is not negotiable is that the
// router validates the same token against the same key: a subgraph
// that trusts a header the router set, without checking it, has
// moved the guard back out of the service again.
builder.Services
    .AddAuthentication()
    .AddJwtBearer(options =>
    {
        options.TokenValidationParameters = new()
        {
            ValidateIssuer = false,
            ValidateAudience = false,
            ValidateLifetime = false,
            IssuerSigningKey = new SymmetricSecurityKey(
                Encoding.UTF8.GetBytes(SigningKey.Value))
        };
    });

builder.Services.AddAuthorization();

builder
    .AddGraphQL()
    // Without this the [Authorize] attributes compile, reach the
    // schema as @authorize, and are never executed.
    .AddAuthorization()
    .AddSpeakersTypes()
    .RegisterDbContextFactory<SpeakerContext>()
    .AddQueryContext()
    .AddPagingArguments()
    .AddGlobalObjectIdentification()
    .AddApolloFederation()
    .ModifyCostOptions(options => options.ApplyCostDefaults = false)
    .TryAddTypeInterceptor<ShareNodeFields>();

var app = builder.Build();

using (var scope = app.Services.CreateScope())
{
    var factory = scope.ServiceProvider
        .GetRequiredService<IDbContextFactory<SpeakerContext>>();
    using var db = factory.CreateDbContext();
    db.Database.EnsureCreated();
}

app.UseAuthentication();
app.UseAuthorization();

app.MapGraphQL()
    .WithOptions(options =>
    {
        // Serve the Nitro build that shipped in the package. The
        // default, ServeMode.Latest, downloads the current one
        // from ChilliCream's CDN instead.
        options.Tool.ServeMode = ServeMode.Embedded;
    });

app.RunWithGraphQLCommands(args);
\end{minted}

\index{authorization!the middleware has to be registered}%
The \code{AddAuthorization()} inside the GraphQL chain is the one to watch.
Without it the attributes below compile, reach the schema as a directive, and
are never executed, which is a schema that documents a rule nobody enforces.
The two \code{app.Use} calls are ASP.NET Core's, and they are what puts a user
on the request before Hot Chocolate looks for one.

\subsection{The Column, and the Guard}

\index{authorization!the column that is not for the program}%
Now the data. \code{Speaker.cs} gains a fourth field:

\begin{minted}[highlightlines={10-16}]{csharp}
namespace Speakers;

/// <summary>
/// A person giving one or more sessions. This service owns the rows;
/// the Sessions subgraph knows only that a speaker with this id exists.
/// </summary>
/// <param name="Id">Stable identifier for this speaker.</param>
/// <param name="Name">The name to print in the program.</param>
/// <param name="Bio">A sentence or two, if the speaker sent any.</param>
/// <param name="Email">Where the program committee writes. Not for
/// the printed program.</param>
public sealed record Speaker(
    int Id,
    string Name,
    string? Bio,
    string Email);
\end{minted}

\index{authorization!the seed the responses come from}%
And \code{SpeakerData.cs} gains three addresses, which are the values every
response in this chapter comes back with:

\begin{minted}[highlightlines={12-21}]{csharp}
namespace Speakers;

/// <summary>
/// The speaker list, seeded into this service's own database. These are
/// the rows chapter 3 seeded into the Sessions service; they live here
/// now, and only here.
/// </summary>
public static class SpeakerData
{
    public static readonly IReadOnlyList<Speaker> Speakers =
    [
        new Speaker(
            1, "Ada Fischer", "Works on query planners.",
            "ada@example.test"),
        new Speaker(
            2, "Bruno Kaminski", null,
            "bruno@example.test"),
        new Speaker(
            3, "Chidi Okafor", "Runs platform engineering.",
            "chidi@example.test")
    ];
}
\end{minted}

\index{authorization!where the attribute goes first}%
The requirement was about the speaker list, and \code{speakers} is the field
that serves the speaker list, so that is where the attribute goes:

\begin{minted}[highlightlines={2,12}]{csharp}
using GreenDonut.Data;
using HotChocolate.Authorization;

namespace Speakers;

[QueryType]
public static class Query
{
    /// <summary>
    /// Everyone speaking at the conference, in program order.
    /// </summary>
    [Authorize]
    [UsePaging(DefaultPageSize = 2, MaxPageSize = 10)]
    public static IQueryable<Speaker> GetSpeakers(
        SpeakerContext db,
        QueryContext<Speaker> query)
        => db.Speakers
            .OrderBy(s => s.Name)
            .With(query);
}
\end{minted}

\subsection{It Works}

\index{authorization!the root field is refused}%
Ask the Speakers service for the list with nothing in the
\code{Authorization} header, on its own port, and it says no:

\begin{minted}[fontsize=\footnotesize,breakanywhere]{text}
POST /graphql HTTP/1.1
Host: localhost:5002
Content-Type: application/json

{"query":"{ speakers(first: 2) { nodes { name email } } }"}
\end{minted}

\begin{minted}[fontsize=\footnotesize,breakanywhere]{json}
{"errors":[
  {"message":"The current user is not authorized to access this resource.",
   "path":["speakers"],
   "extensions":{"code":"AUTH_NOT_AUTHENTICATED"}}],
 "data":{"speakers":null}}
\end{minted}

\index{statement count!a refusal before the resolver}%
The refusal is clean and it is cheap. The path names the field, the extension
gives a code a client can branch on, and
chapter~\ref{ch:data-without-n-plus-1}'s interceptor prints nothing at all:
the middleware runs in front of the resolver, so no statement reaches SQLite.
Nothing was read in order to be refused.

\subsection{And Then It Does Not}

\index{entities@\code{\_entities}!answers the same rows unguarded}%
Here is the same service, the same process, the same missing header, one
field along:

\begin{minted}[fontsize=\footnotesize,breakanywhere]{graphql}
query E($r: [_Any!]!) {
  _entities(representations: $r) { ... on Speaker { id name email } }
}
\end{minted}

\begin{minted}[fontsize=\footnotesize,breakanywhere]{json}
{"r":[{"__typename":"Speaker","id":"U3BlYWtlcjox"}]}
\end{minted}

\begin{minted}[fontsize=\footnotesize,breakanywhere]{json}
{"data":{"_entities":[{"id":"U3BlYWtlcjox","name":"Ada Fischer",
 "email":"ada@example.test"}]}}
\end{minted}

\index{entities@\code{\_entities}!no error key}%
There is no error key and nothing for a client to branch on. The log shows
one statement, the same statement the reference resolver has issued since
chapter~\ref{ch:second-third-subgraph}, and the column the attribute was
written to protect is in the response.

\index{entities@\code{\_entities}!the whole table in one request}%
Nor is it one row at a time. The entity route takes a list, which is the
whole reason chapter~\ref{ch:router-n-plus-1} could batch through it, and a
list is what it answers:

\begin{minted}[fontsize=\footnotesize,breakanywhere]{json}
{"r":[{"__typename":"Speaker","id":"U3BlYWtlcjoz"},
      {"__typename":"Speaker","id":"U3BlYWtlcjox"},
      {"__typename":"Speaker","id":"U3BlYWtlcjoy"}]}
\end{minted}

\begin{minted}[fontsize=\footnotesize,breakanywhere]{json}
{"data":{"_entities":[
  {"id":"U3BlYWtlcjoz","name":"Chidi Okafor",
   "email":"chidi@example.test"},
  {"id":"U3BlYWtlcjox","name":"Ada Fischer",
   "email":"ada@example.test"},
  {"id":"U3BlYWtlcjoy","name":"Bruno Kaminski",
   "email":"bruno@example.test"}]}}
\end{minted}

\index{node@\code{node(id:)}!the third door}%
And there is a third door, which chapter~\ref{ch:schema-design} built and
nobody thought of as a door. \code{node(id:)} takes a global node id and
returns whatever it identifies:

\begin{minted}[fontsize=\footnotesize,breakanywhere]{graphql}
{ node(id: "U3BlYWtlcjox") { ... on Speaker { name email } } }
\end{minted}

\begin{minted}[fontsize=\footnotesize,breakanywhere]{json}
{"data":{"node":{"name":"Ada Fischer","email":"ada@example.test"}}}
\end{minted}

Figure~\ref{fig:ch12-three-doors} is the shape of it.

\begin{figure}[htbp]
  \centering
  % Three doors into one column of data, and a guard standing in front of
% exactly one of them.
%
% The Speakers subgraph exposes three fields on Query, all three advertised
% by introspection, and all three reach the same rows of the Speaker table:
%   speakers(first:)                the root field
%   _entities(representations:)     the entity route the router uses
%   node(id:)                       the node field
% A guard placed on the root field only stops speakers(first:) and does
% nothing to the other two: _entities and node still return Speaker.email in
% full. That is the whole picture: one path blocked, two paths open, one
% column of data behind all three.
%
% Layout is boxes and links, the idiom figures/tikz/ch11-position-is-the-join.tex
% uses, not the axis-and-ticks timeline of ch01-one-field.tex and
% ch06-two-routes.tex: the subject here is three parallel routes converging
% on one target, not a sequence unfolding over one request, so there is no
% single shared axis for all three to sit on.
%
% Three field boxes on the left, one tall box on the right standing for the
% Speaker rows and their email column (drawn tall on purpose, so a single
% box can receive an arrow at more than one height and still read as one
% target). The entities and node links run straight through from field box
% to data box, arrowhead at the data box: the request is answered. The
% speakers link runs only as far as a small guard box and stops there: it
% never reaches the data box at all, which is the point a reader should
% get in one glance, not from a label.
%
% A short dashed mark crosses the entities and node links at the same
% x-coordinate as the guard box, purely as an annotation (the links
% themselves are not broken there). It marks where a check would also have
% to sit if it were checking the data rather than one field, and the note
% beside it says so. This is the one place this figure borrows the "guard
% belongs around the data, not on one path" idea the brief allows; it stays
% inside one panel rather than adding a second, because the same guard box
% and the same gate x-coordinate already carry the comparison.
%
% No fills, square corners, one stroke weight throughout (0.4pt), arrowhead
% only where a request actually lands (the data box, or the guard box for
% the one path that gets no further) -- the same "arrowhead marks the one
% directional claim" idiom ch11 uses. Dashed is a short, local mark here,
% not a link of its own, and does not carry the same meaning dashed carries
% in ch11; the note beside it is what supplies the meaning, same as ch11's
% own legend does for its two link styles.
%
% No versions, timings or counts anywhere, per SPEC decision 19.
%
% No quotation marks anywhere in this file (strict Ascii mode; a literal "
% reads as a quote character to the prose gate, and \" is a diaeresis
% accent that fails to compile inside a node). \_entities is written with
% the escape TikZ needs for an underscore in text mode, the same way
% ch06-two-routes.tex escapes \_\_typename.
%
% Widths: field boxes are 38mm text plus 2mm inner sep, centered at x=22,
% so they span x=1 to x=43, starting near the left margin the way every
% other figure in this book does. The guard box is centered at x=84 and
% spans x=78 to x=90. The data box is 26mm text plus 2mm inner sep,
% centered at x=135, spanning x=120 to x=150 -- the same right-hand edge
% ch11's note uses, and inside the 155mm text measure (A4, 30mm inner,
% 25mm outer). The note beside the dashed marks runs from x=90 to x=136,
% clear of that edge.
%
% Bare tikzpicture (house style): the figure environment, caption and
% label live at the call site in the section file.
\begin{tikzpicture}[
  x=1mm, y=1mm,
  every node/.append style={font=\sffamily\scriptsize},
  lane/.style={anchor=west, font=\sffamily\scriptsize\itshape},
  fieldbox/.style={draw, line width=0.4pt, align=center, text width=38mm,
    inner sep=2mm, minimum height=14mm},
  gatebox/.style={draw, line width=0.4pt, align=center, text width=10mm,
    inner sep=1mm, minimum height=10mm},
  databox/.style={draw, line width=0.4pt, align=center, text width=26mm,
    inner sep=2mm, minimum height=100mm},
  reachlink/.style={draw, line width=0.4pt,
    -{Straight Barb[length=1.4mm, width=1.6mm]}},
  blocklink/.style={draw, line width=0.4pt,
    -{Straight Barb[length=1.4mm, width=1.6mm]}},
  gapmark/.style={draw, line width=0.4pt, dashed},
  note/.style={align=left, inner sep=2pt},
]

% ---------------------------------------------------------------- heading
\node[lane] at (0,58) {Speakers subgraph, three fields of Query};

% ------------------------------------------------------------ field boxes
\node[fieldbox] (fspk)  at (22,40)  {speakers(first:)\\the root field};
\node[fieldbox] (fent)  at (22,0)
  {\_entities(representations:)\\the entity route};
\node[fieldbox] (fnode) at (22,-40) {node(id:)\\the node field};

% -------------------------------------------------------------- the guard
\node[gatebox] (guard) at (84,40) {guard};
\node[note, anchor=north] at (84,32) {refused here};

% ------------------------------------------------------------- data column
\node[databox] (data) at (135,0) {Speaker rows\\email column};

% ------------------------------------------------------------------ links
% The two open routes: unbroken, arrowhead at the data box.
\draw[reachlink] (fent.east)  -- (120,0);
\draw[reachlink] (fnode.east) -- (120,-40);

% The guarded route: stops at the guard box and goes no further.
\draw[blocklink] (fspk.east) -- (guard.west);

% ------------------------------------------------------- where a check on
% ------------------------------------------------- the data would also sit
\draw[gapmark] (84,-5)  -- (84,5);
\draw[gapmark] (84,-45) -- (84,-35);
% text width 26mm, not 46: the note is anchored at the gate x and the data
% box starts at x=120, so anything wider than 30mm runs into the box.
\node[note, anchor=west, text width=26mm] at (90,-20)
  {the same crossing point on both open routes, and nothing stands there};

\end{tikzpicture}

  \caption{One column of data and three fields of \code{Query} that reach it.
  The attribute went on the route the requirement named, and the requirement
  was written in terms of a screen rather than in terms of a column. The two
  open routes are drawn unbroken because nothing interrupts them: the dashed
  marks are not breaks but the place a check would have to stand to be
  checking the data instead of one field, and nothing stands there. A guard
  that has to be repeated once per field is a guard that will be forgotten
  the next time somebody adds one.}
  \label{fig:ch12-three-doors}
\end{figure}

\subsection{None of This Has to Be Guessed}

\index{introspection!the entity route is advertised}%
The awkward part is not that the entity route exists. It is that the service
will describe it to anybody who asks. Introspect \code{Query} with no
credentials:

\begin{minted}[fontsize=\footnotesize,breakanywhere]{json}
{"data":{"__type":{"fields":[{"name":"node"},{"name":"nodes"},
 {"name":"speakers"},{"name":"_service"},{"name":"_entities"}]}}}
\end{minted}

\index{introspection!the entity union names the types}%
Then ask what the route can return, because the schema says that too:

\begin{minted}[fontsize=\footnotesize,breakanywhere]{json}
{"data":{"__type":{"kind":"UNION","possibleTypes":[{"name":"Speaker"}]}}}
\end{minted}

\index{service@\code{\_service}!hands over the schema}%
The one thing left to work out is the shape of a representation, and
\code{\_service} settles it. Chapter~\ref{ch:first-subgraph} used that field
to read the schema this service exports. So can anyone else, without a token,
including the \code{@key} that says which field a representation carries.

\index{authorization!the guard is on a door, not on the data}%
Three fields, one column, one guard on one of them. The attribute is not
wrong and it is not broken. It is in the wrong place, and the wrong place is
the one the requirement pointed at.
\index{authorization!the guard on the root field|)}%
